\section*{第 \ref{ch:b3:palaeomagnetism} 章——地磁学与古地磁学}

\begin{solution}{exo:b3:palaeomagnetism:1}
$\tan I = 2\tan\ang{30} = 1.1547$，故 $I = \ang{49.1}$。反之，$\tan\lambda
= \tfrac12\tan(\ang{-20}) = -0.1820$，故 $\lambda = \ang{-10.3}$，即
南纬 \ang{10.3}。
\end{solution}

\begin{solution}{exo:b3:palaeomagnetism:2}
$H = \sqrt{20000^2 + 2000^2} = \qty{20100}{nT}$；
$F = \sqrt{20000^2 + 2000^2 + 40000^2} = \qty{44770}{nT}$。磁偏角
为 $\arctan(-2000/20000) = \ang{-5.7}$，即北偏西 \ang{5.7}；
\omterm{def:b3:palaeomagnetism:inclination}{磁倾角}为 $\arctan(40000/20100) = \ang{63.3}$，指向下方。
\end{solution}

\begin{solution}{exo:b3:palaeomagnetism:3}
赤道处 $B_0 = -g_1^0 = \qty{29350}{nT}$，
两极处 $2B_0 = \qty{58700}{nT}$。$m = 4\pi a^3B_0/\mu_0 = a^3B_0/\num{e-7} = \num{2.5862e20}\times
\num{2.935e-5}/\num{e-7} = \qty{7.59e22}{A.m^2}$，略小于
全偶极子的 \qty{7.69e22}{A.m^2}，后者还包含 $g_1^1$ 和 $h_1^1$。
\end{solution}

\begin{solution}{exo:b3:palaeomagnetism:4}
(a) \omterm{def:b3:palaeomagnetism:nrm}{热剩磁}，记录了熔岩冷却时
经过其颗粒\omterm{def:b3:palaeomagnetism:blocking}{阻挡温度}时的磁场。
(b) \omterm{def:b3:palaeomagnetism:nrm}{碎屑剩磁}，记录了泥沙沉降
并因压实而被固定期间的磁场。
(c) \omterm{def:b3:palaeomagnetism:nrm}{化学剩磁}，记录了赤铁矿生长
期间的磁场，即在早期埋藏时形成，因而略晚于沉积。
(d) 沿现今磁场方向的粘滞剩磁，由弛豫时间最短的颗粒叠加到原始\omterm{def:b3:palaeomagnetism:nrm}{热剩磁}上；
它在退磁的最初几步即被去除。
\end{solution}

\begin{solution}{exo:b3:palaeomagnetism:5}
单位矢量 $(\cos I\cos D,\ \cos I\sin D,\ \sin I)$ 为
\begin{gather*}
  (0.6545, 0.1391, 0.7431),\quad (0.6670, 0.2695, 0.6947),\\
  (0.5235, 0.1501, 0.8387),\quad (0.5985, 0.0629, 0.7986).
\end{gather*}
其和为 $(2.4435, 0.6216, 3.0751)$，长度 $R = 3.9766$。平均值：
$D = \arctan(0.6216/2.4435) = \ang{14.3}$，$I = \arcsin(3.0751/3.9766) =
\ang{50.7}$。$k = 3/(4 - 3.9766) = 128$。取 $20^{1/3} - 1 = 1.7144$，
$\cos\alpha_{95} = 1 - (0.0234/3.9766)\times1.7144 = 0.98991$，故
$\alpha_{95} = \ang{8.1}$：四个样本时，锥角始终较宽，
无论该组数据多么集中。
\end{solution}

\begin{solution}{exo:b3:palaeomagnetism:6}
$\lambda = \arctan(\tfrac12\tan\ang{50}) = \ang{30.79}$，故 $p = \ang{59.21}$。
\[\sin\lambda_p = \sin\ang{40}\cos\ang{59.21} + \cos\ang{40}\sin\ang{59.21}
\cos\ang{20} = 0.9474,\] 故 $\lambda_p = \ang{71.34}$（北纬）。$\sin\beta =
\sin\ang{59.21}\sin\ang{20}/\cos\ang{71.34} = 0.9184$，$\beta = \ang{66.66}$。
由于 $\cos p = 0.5119 < \sin\ang{40}\times0.9474 = 0.6090$，
故适用第二种情形：$\varphi_p = 10 + 180 - 66.66 = \ang{123.34}$（东经）。极点位于
北纬 \ang{71.3}，东经 \ang{123.3}。
\end{solution}

\begin{solution}{exo:b3:palaeomagnetism:7}
(a) $(7.69 - 8.32)/8.32 = \qty{-7.6}{\percent}$，历时 125 年，
即每世纪 \qty{-6.1}{\percent}。(b) $29350/12.6 = 2329$，约 2300 年。
这并非预报：变化率本身随五年一档的模型不同而改变，
过去偶极矩曾经既更高又更低，却从未发生倒转，
而磁场倒转的时间间隔不规则，
任何变化率都无法预测。
\end{solution}

\begin{solution}{exo:b3:palaeomagnetism:8}
在 20 至 \qty{300}{\celsius} 之间，矢量损失了 $(6.4, 0.0, 13.6)$：
一个分量为 \qty{15.0}{mA/m}，$D = \ang{0.0}$ 且
$I = \arctan(13.6/6.4) = \ang{64.8}$。超过 \qty{300}{\celsius} 后，矢量损失了 $(25.8, 18.1, 24.6)$：
\qty{40.0}{mA/m}，$D = \arctan(18.1/25.8) = \ang{35.1}$，$I =
\arctan(24.6/31.5) = \ang{38.0}$。第二个分量指向原点，是
\omterm{def:b3:palaeomagnetism:zijderveld}{特征剩磁} ChRM；第一个分量沿北向陡倾方向，在
低温下被去除，是沿现今磁场方向的粘滞叠加。
\end{solution}

\begin{solution}{exo:b3:palaeomagnetism:9}
$\lambda = \arctan(\tfrac12\tan(\ang{-35})) = \ang{-19.3}$。对于 $I = \ang{-40}$
和 \ang{-30}：\ang{-22.8} 和 \ang{-16.1}，区间长为 \ang{6.7}。
$p = \ang{90} - \lambda = \ang{109.3}$，$\cos^2p = 0.1096$：
$\mathrm dp = 2.5\times(1 + 0.329) = \ang{3.3}$，且 $\mathrm dm =
5\times\sin\ang{109.3}/\cos\ang{35} = \ang{5.8}$。
\end{solution}

\begin{solution}{exo:b3:palaeomagnetism:10}
(a) $\tan\lambda = \tfrac12$，$\lambda = \ang{26.6}$。$F$ 在赤道处为 $B_0$，
两极处为 $2B_0$：比值为 2。(b) 由 $\tan I = 2\tan\lambda$
可得 $(1 + \tan^2I)\,\mathrm dI = 2(1 + \tan^2\lambda)\,\mathrm d\lambda$，于是
\[
  \frac{\mathrm dI}{\mathrm d\lambda} = \frac{2/\cos^2\lambda}{1 + 4\tan^2\lambda}
  = \frac{2}{\cos^2\lambda + 4\sin^2\lambda} = \frac{2}{1 + 3\sin^2\lambda}.
\]
误差 $\delta I$ 对应 $\delta\lambda = \tfrac12(1 + 3\sin^2\lambda)\,\delta I$，
从赤道处的 $\tfrac12\delta I$ 变为两极处的 $2\,\delta I$：
最大误差出现在极地附近。(c) 在赤道附近，\omterm{def:b3:palaeomagnetism:inclination}{磁倾角}随纬度的变化率
是纬度变化速率的两倍，因此同样的\omterm{def:b3:palaeomagnetism:inclination}{磁倾角}误差
所对应的纬度误差只有一半；在极地附近，\omterm{def:b3:palaeomagnetism:inclination}{磁倾角}几乎不随纬度变化。
\end{solution}

\begin{solution}{exo:b3:palaeomagnetism:11}
(a) $\cos p = \sin(\ang{-35})\sin\ang{70} + \cos\ang{35}\cos\ang{70}\cos\ang{150}
= -0.5390 - 0.2426 = -0.7816$，$p = \ang{141.41}$。于是 $\tan I_x = 2\cot p =
-2.506$，$I_x = \ang{-68.2}$，且 $\lambda_x = \ang{90} - p = \ang{-51.4}$。
$\cos D_x = (\sin\ang{70} - \sin(\ang{-35})\cos p)/(\cos\ang{35}\sin p) =
0.9617$，由于极点位于测站子午线以东（$\sin\ang{150} >
0$），故 $D_x = \ang{15.9}$。(b) $\lambda_o = \arctan(\tfrac12\tan(\ang{-60})) =
\ang{-40.9}$。$R = 30.0 - 15.9 = \ang{14.1}$：岩块
绕竖直轴发生了顺时针旋转。$P = \lambda_x - \lambda_o = -51.4 + 40.9 =
\ang{-10.5}$：该地层位于比稳定板块所预期的位置偏北 \ang{10.5} 处，
故岩块相对板块已向南移动了约 \ang{10.5}，
即约 \qty{1170}{km}。
\end{solution}

\begin{solution}{exo:b3:palaeomagnetism:12}
(a) 取 $\mathrm d\Omega = 2\pi\sin\theta'\,\mathrm d\theta'$，
\[
  \int_0^\theta\frac{\kappa e^{\kappa\cos\theta'}}{4\pi\sinh\kappa}\,2\pi\sin\theta'\,
  \mathrm d\theta' = \frac{e^\kappa - e^{\kappa\cos\theta}}{2\sinh\kappa}
  = \frac{e^\kappa - e^{\kappa\cos\theta}}{e^\kappa - e^{-\kappa}} .
\]
(b) 当 $\kappa$ 较大时，$e^{-\kappa}$ 可忽略；两边除以 $e^\kappa$
余下 $1 - e^{-\kappa(1 - \cos\theta)} \approx 1 - e^{-\kappa\theta^2/2}$。该值
等于 $1 - e^{-1} = \qty{63}{\percent}$，此时 $\kappa\theta^2/2 = 1$，$\theta =
\sqrt{2/\kappa}$ 弧度，或 $\ang{81}/\sqrt\kappa$。(c) $81/\sqrt{50} =
\ang{11.5}$；精确公式给出 \qty{63.1}{\percent}。
\end{solution}

\begin{solution}{pb:b3:palaeomagnetism:1}
\textbf{1.} 红层中的赤铁矿可能由氧化铁和氢氧化铁形成，
且常在沉积物沉积之后很久才生成，从而形成
较晚年代的\omterm{def:b3:palaeomagnetism:nrm}{化学剩磁}。若岩层在沉积后不久即发生褶皱，则正向褶皱检验，
或对两种极性兼有、间距很近的岩层作正向倒转检验，
均可表明该剩磁形成于较早时期。

\textbf{2.} $(\cos\ang{37}\cos\ang{50}, \cos\ang{37}\sin\ang{50}, \sin\ang{37})
= (0.5134, 0.6118, 0.6018)$。

\textbf{3.} $R = \sqrt{3.8943^2 + 2.6500^2 + 3.5717^2} = 5.9114$。$D =
\arctan(2.6500/3.8943) = \ang{34.2}$，$I = \arcsin(3.5717/5.9114) = \ang{37.2}$。

\textbf{4.} $k = 5/(6 - 5.9114) = 56.5$。取 $20^{1/5} - 1 = 0.8206$，
$\cos\alpha_{95} = 1 - (0.0886/5.9114)\times0.8206 = 0.98771$，$\alpha_{95} =
\ang{9.0}$。

\textbf{5.} B：$k = 6/(7 - 6.8979) = 58.8$；$20^{1/6} - 1 = 0.6475$，
$\cos\alpha_{95} = 1 - (0.1021/6.8979)\times0.6475 = 0.99042$，$\alpha_{95} =
\ang{7.9}$。C：$k = 7/(8 - 7.8707) = 54.1$；$20^{1/7} - 1 = 0.5340$，
$\cos\alpha_{95} = 1 - (0.1293/7.8707)\times0.5340 = 0.99123$，$\alpha_{95} =
\ang{7.6}$。

\textbf{6.} B 的 $k$ 最大，分组最为紧密；C 的
$\alpha_{95}$ 最小，平均值最为可靠，因其采样点最多。三者
相差不大：三个地层的分组都很好。

\textbf{7.} 三者均为北纬：$\lambda_A = \arctan(\tfrac12\tan\ang{37.17}) = \ang{20.76}$，
$\lambda_B = \ang{32.15}$，$\lambda_C = \ang{45.33}$。

\textbf{8.} $p = 90 - 20.76 = \ang{69.24}$。$\sin\lambda_p =
\sin\ang{47}\cos\ang{69.24} + \cos\ang{47}\sin\ang{69.24}\cos\ang{34.23} =
0.2593 + 0.5272 = 0.7865$，故 $\lambda_p = \ang{51.86}$ 北纬。

\textbf{9.} $\sin\beta = \sin\ang{69.24}\sin\ang{34.23}/\cos\ang{51.86} =
0.8516$，$\beta = \ang{58.40}$。由于 $\cos p = 0.3545 < \sin\ang{47}\times
0.7865 = 0.5752$，故适用第二种情形，$\cos\beta < 0$：
$\varphi_p = 15.00 + 180 - 58.40 = \ang{136.60}$ 东经。

\textbf{10.} $\cos^2p = 0.1257$：$\mathrm dp = 4.50\times1.377 = \ang{6.2}$；
$\mathrm dm = 8.99\times\sin\ang{69.24}/\cos\ang{37.17} = \ang{10.5}$。

\textbf{11.} 由 \qty{60}{Ma} 时的北纬 \ang{51.9} 移至北纬 \ang{66.5}
（\qty{30}{Ma}）及北纬 \ang{87.1}（\qty{4}{Ma}），三者近乎位于同一经线上
（东经 \ang{136}--\ang{139}）：一条指向地理极的路径，
最年轻的地层已落入其锥形范围之内。

\textbf{12.} 火山灰层位于 \qty{22}{m} 处，属第二个正极性带，
19.7--\qty{24.6}{m}；年龄约为 \qty{3.2}{Ma}，该带为 C2An.2n
（3.207--3.116）。随后 4.2--\qty{14.8}{m} 带为 C2An.3n，
29.2--\qty{54.0}{m} 带为 C2An.1n，反极性底界属于 C2An 之下的极性时，
C2An，而反极性顶界属于其上的极性时。

\textbf{13.} $(19.7 - 4.2)/(3.596 - 3.207) = 15.5/0.389 = \qty{39.8}{m/Ma}$；
$(54.0 - 24.6)/(3.116 - 2.581) = 29.4/0.535 = \qty{55.0}{m/Ma}$。

\textbf{14.} 底界：$3.596 + 4.2/39.8 = \qty{3.70}{Ma}$。顶界：$2.581 -
10.0/55.0 = \qty{2.40}{Ma}$。

\textbf{15.} 仅对倒转界线定年，界线之间相隔约几十万年；
其间假定速率恒定。沉降、
沉积物供给或气候的变化，或沉积间断，均会改变速率，而
间断甚至可能删去一个短的带，致使对比产生误导。

\textbf{16.} $\cos p = \sin\ang{47}\sin\ang{78} + \cos\ang{47}\cos\ang{78}
\cos\ang{175} = 0.7154 - 0.1413 = 0.5741$，$p = \ang{54.96}$。$\tan I_x =
2\cot p = 1.4023$，$I_x = \ang{54.5}$；$\lambda_x = \ang{35.04}$ 北纬。

\textbf{17.} $\cos D_x = (\sin\ang{78} - \sin\ang{47}\times0.5741)/
(\cos\ang{47}\sin\ang{54.96}) = 0.99976$；极点位于测站
经线以东（$\sin\ang{175} > 0$），故 $D_x = \ang{1.3}$。

\textbf{18.} 由 $I_x$ = \ang{54.5}、\ang{60.0} 及
\ang{63.6}时，预期\omterm{def:b3:palaeomagnetism:vgp}{古纬度}分别为 \ang{35.03}、\ang{40.89} 及 \ang{45.21}。$P_A = 35.03 - 20.76 =
\ang{14.3}$，$P_B = 40.89 - 32.15 = \ang{8.7}$，$P_C = 45.21 - 45.33 =
\ang{-0.1}$。

\textbf{19.} $R_A = 34.2 - 1.3 = \ang{32.9}$，$R_B = 23.9 - 2.9 = \ang{21.0}$，
$R_C = 3.4 - 2.5 = \ang{0.9}$，均为顺时针。

\textbf{20.} $\Delta D_o = \arcsin(\sin\ang{8.99}/\cos\ang{37.17}) = \ang{11.3}$；
$\Delta D_x = \arcsin(\sin\ang{3.0}/\cos\ang{54.5}) = \ang{5.2}$；$\Delta R_A =
\sqrt{11.3^2 + 5.2^2} = \ang{12.4}$。

\textbf{21.} $R_A = \ang{32.9} \pm \ang{12.4}$ 与 $R_B = \ang{21.0} \pm
\ang{13.8}$ 均显著；$R_C = \ang{0.9} \pm \ang{17.9}$ 则不显著。

\textbf{22.} 6000 万年前，凯斯塔的纬度比稳定大陆所确定的位置偏南约 \ang{14}，
且其方位较现今逆时针偏转约 \ang{33}。
至 \qty{30}{Ma} 时，它已北移约 \ang{14} 中的
\ang{6}，并顺时针旋转了 \ang{33} 中的 \ang{12}；
至 \qty{4}{Ma} 时，它已到达大陆上的现今位置，
此时已无剩余旋转。它沿大陆边缘向北漂移，
一边漂移一边顺时针转动，并于 \qty{4}{Ma} 之前已就位。

\textbf{23.} $47.00 - 20.76 = \ang{26.24}$，或 $26.24\times111.19 =
\qty{2918}{km}$。此数包括稳定大陆本身的运动；
相对于它，凯斯塔移动了 $P_A = \ang{14.3}$。

\textbf{24.} $I = \ang{28.18}$ 和 \ang{46.16} 分别给出 $\lambda = \ang{15.00}$ 和
\ang{27.50}：半宽为 \ang{6.25}，一倍标准差为
$6.25/1.96 = \ang{3.19}$，即 \qty{355}{km}。自地层 A 形成以来，凯斯塔已
向北漂移了
\textbf{\qty[text-series-to-math=true]{2918 \pm 355}{km}}。
\end{solution}
